\chapter*{Аннотация}
\addcontentsline{toc}{chapter}{Аннотация}

Работа посвящена исследованию подходов к формированию групповых музыкальных рекомендаций и оценке возможностей их практического применения в сервисах музыкального сопровождения заведений и мероприятий. 

Разработаны теоретический концепт сервиса на базе инфраструктуры «Яндекс Музыки», практическая реализация алгоритма групповых рекомендаций на базе открытого датасета YAMBDA (Yandex Music Billion-Interactions Dataset), а также механизм фиксации нахождения пользователей в заведении с помощью технологии BLE (Bluetooth Low Energy).

Работа логически разделена на бизнес-часть (главы 1-2), техническую часть (главы 3-5) и описание демонстрационных прототипов (глава 6).

В бизнес-части проанализированы рынок музыки для бизнеса, конкуренты, правовые аспекты публичного воспроизведения, целевые сегменты и модель монетизации. Предложена концепция сервиса «Яндекс Музыка для бизнеса» и функции «Резонанс», адаптирующей звучащую в заведении музыку под посетителей при сохранении приватности и управляемости со стороны бизнеса.

Техническая часть построена как двухуровневая схема: замороженный персональный последовательный скорер plain SASRec с BCE-функцией потерь и обучаемые групповые агрегаторы поверх кэша персональных скоров. Скорер воспроизводит индустриальный SASRec-бейзлайн на YAMBDA: test NDCG@10 = 0,0726 на Listen+ против 0,0742. Для синтетических эфемерных групп размера 2-5 сравниваются ID-based AGREE, GroupIM, Audio-AGREE и Group Cross-Attention. Показано, что аудиоагрегаторы значимо превосходят ID-based бейзлайны по NDCG@10 и NDCG@20, особенно для групп среднего размера.

В заключительной части описаны демонстрационные прототипы: визуализация работы сервиса и модели рекомендаций на схеме помещения с анонимизированными пользователями и треками из YAMBDA, а также Android-приложение, иллюстрирующее определение нахождения посетителя в заведении по силе сигнала BLE-маяка. Также рассмотрены потенциальные направления развития сервиса и технологии.

\section*{Ключевые слова}

Групповые музыкальные рекомендации; коллективные рекомендации; музыкальные рекомендательные системы; последовательные рекомендации; SASRec; AGREE; GroupIM; Audio-AGREE; Group Cross-Attention; аудиоагрегаторы; YAMBDA; Yandex Music Billion-Interactions Dataset; BLE; Bluetooth Low Energy; музыка для бизнеса; публичное воспроизведение музыки; РАО; ВОИС; Яндекс Музыка.
